\section{硅碳共生的文明宪法}

个体与组织各自能做的事已经交代完，它们都需要一个更大的东西来承载。组织嵌在文明里，组织的评价函数受制于它所在的文明关心什么。把前面四个原则放在一起看，它们更接近一部宪法的条款而不是一份技术方案：宪法不规定做什么，它规定当不同的价值发生冲突时按什么次序裁决。这一节把尺度放大到文明，并把前面的原则收成一组可以被反复引用的约束。

\subsection{文明的高阻力前沿}

\subhead{从组织到文明}

上一节说明可信组织如何维护慢变量，这里把尺度再放大一层。组织是嵌在更大的系统里的，它的评价函数受制于它所在的文明关心什么。如果整个文明把资源持续投向反馈最快、风险最低、回报最确定的方向，那么任何单个组织维持慢变量的努力都会被慢慢耗尽。因此需要问一个更大的问题：文明的资源应当投向哪里，以及这个问题的答案在生成能力变得充足之后有没有变化。

这个问题容易被两种话语占据。一种是技术乐观的扩张叙事，把深空、聚变、生命工程列成清单，默认只要能力足够就应该去做；另一种是内敛的收缩叙事，主张把资源用于解决当下的具体问题，认为向外的投入是浪费。这两种叙事都没有触及本书关心的地方，因为它们争论的是投多少，而不是这些方向有什么共同的结构性质。

\subhead{高阻力前沿的共同结构}

把这些方向放在一起看，它们的共同点不是宏大，而是它们都同时具备三样约束。第一是资源势垒高。深空探测要面对推进剂、辐射防护、通信延迟和载荷质量之间的硬性交换，深海作业要面对压力、腐蚀和能源供给，生态修复要面对土壤、水文、物种关系和几十年尺度的演替，基础科学设施要面对材料、真空、低温、精密制造和巨额前期投入。这些约束不能通过更好的算法绕过，因为它们来自物理定律和材料的实际极限。

第二是时间刚性问题突出。航天任务的窗口由轨道力学决定，深海作业受制于海况和装备寿命，生态系统的恢复有自己的演替节奏，材料的老化和生物的发育都有内禀速率，基础科学的验证往往需要下一代仪器。这些过程的时间尺度不由人的意愿决定，压缩它们要么不可能，要么会改变过程的性质。第三是拓扑耦合密集。这些方向上的任何一步都牵动大量互相依赖的子系统，且各子系统的优化目标不一致：一个更轻的结构可能降低安全裕度，一个更快的迭代可能牺牲验证深度，一个更省钱的方案可能把风险转给后期维护。耦合越密集，局部优化的收益越容易被系统层面的代价抵消。

三样约束同时存在的方向，就是这里所说的高阻力前沿。写成集合，
\begin{equation}
\mathcal{F}=\bigl\{(R,T,G)\ :\ R\ge R_{*},\ T\ge T_{*},\ G\ge G_{*}\bigr\}
\label{eq:frontier}
\end{equation}
式中 $R$、$T$、$G$ 仍是资源势垒、时间刚性与拓扑耦合三个坐标，与前面的定义同一套刻度；$R_{*}$、$T_{*}$、$G_{*}$ 分别是三条门槛，即“高”“突出”“密集”各自的临界取值，三者的具体位置随领域而定，本书不给出数值。这个式子里三个条件是合取而不是析取，这是它唯一的要点：只满足其中一条或两条的方向不属于 $\mathcal{F}$。一条方向若仅仅资金门槛高，而过程可以在几年内完成、牵动的子系统也不多，那它是昂贵，不是高阻力；一条方向若仅仅牵连面广，而投入和时间都不苛刻，那它是复杂，也不是高阻力。深空探测、生态修复与基础科学之所以能被放在一起讨论，正是因为它们同时越过三条门槛，而不是因为它们各自在某一维上特别突出。这个定义不涉及价值，只涉及结构，因此它可以把看起来完全不同的领域放在一起讨论。

\subhead{为什么这些方向重要}

高阻力前沿的价值来自值得获得的知识、服务或未来可能性。它们要求方案经受实际条件检验；这些检验也可能昂贵、缓慢而稀疏。需要投入的是有区分力的试验、阶段性承诺、恢复路径与完整记录，让有限接触产生足够信息。

更具体地说，这些方向迫使参与者重新接触几样东西。一是材料。材料的性质不可协商，它不因为被描述得好而改变。二是时间。任务的周期、发育的周期、演替的周期都不能被压缩，参与者必须学会在长周期下工作。三是风险。这些方向上的失败是有代价的，而且代价无法完全转移给保险或第三方。四是协作。这些系统的复杂度超过任何单个小组的处理能力，因此必须建立跨机构、跨学科、跨代际的协作结构。五是承责。一次失败必须有人承担，而承担责任的人必须理解自己在系统中的位置。

这五样东西在低黏滞环境中都可以被削弱而不产生即时后果，这正是它们容易被忽略的原因。高阻力前沿的作用是把它们重新变成不可回避的条件。

\subhead{不是对困难的崇拜}

把高阻力前沿说成重要，容易被读成对困难的崇拜，这个读法需要被明确排除。困难本身没有价值，这一点在本书前面已经说过一次，在这里需要再说一次，因为这一节离浪漫化只有一步之遥。为了让论证站得住，可以给出一个反向检验：如果一个方向的阻力只是无效摩擦（流程繁琐，审批冗长，工具落后），那么它不构成高阻力前沿，它只是被浪费的资源。判定标准仍然是摩擦的功能，不是摩擦的大小。

另一条需要排除的读法是把高阻力前沿等同于星辰大海。向外的扩张不是唯一的选项，也不是内定的方向。生态修复、基础科学、公共教育制度的重建同样符合前面的三个条件，而且它们的耦合密度往往更高，因为它们面对的是活的和人的系统。把高阻力前沿等同于航天是一种流行但并不准确的联想，它来自这些领域在公众视野中的可见度，而不是来自它们的结构性质。

还有第三条限制需要说明，这条限制针对的是投入规模。高阻力前沿的回报周期通常在几十年以上，而它们要求的投入是不可逆的。这意味着投向这些方向的资源在其他地方不会被投入，因此每一笔投入都需要一个价值判断，而不是一个技术判断。什么人有权决定把几十年的资源押在一个方向上，以及当前的人凭什么替未来的人做这个决定，这两个问题在前面关于目标设定的讨论里出现过，在这里重新出现，而且更尖锐，因为这里的时间尺度超出了任何一代人的任期。

\subhead{前沿的另一个功能}

除了作为能力和校准的来源，高阻力前沿还有一个功能值得单独指出：它是低黏滞系统无法完全吸收的资源池。当一个文明的所有活动都可以在虚拟空间里以低代价完成时，那些不能被虚拟化的需求会失去供给，因为没有组织愿意去做。这些需求不会消失，它们会在某个时刻以危机的形式出现：基础设施老化到无法继续维护，生态退化到无法恢复，基础科学长期没有新的实验证据，教育系统培养出来的能力与实际问题脱节。

从投入产出的角度看，预防这些危机的支出永远不划算，因为它们的收益是“没有发生的灾难”。这个不对称在前面的内容里反复出现，它解释了为什么高阻力前沿需要制度性的保护而不能靠市场自发供给。不同之处在于这里的尺度更大：保护一个实验室的长期方向可以由一个基金会完成，保护一个文明的高阻力前沿需要的是持续几十年的公共承诺，而公共承诺的期限通常短于这些方向的回报周期。

这个错配是长期的，没有简单的解法。可以做的至少有一件事：把高阻力前沿的存在理由说清楚，使它不被当作少数人的偏好。它们不是文明的点缀，而是文明保留自己与真实世界接触的接口的地方。如果这些接口被关闭，剩下的活动仍然可以繁荣，但它们繁荣的是一个封闭的系统，其中的判断不再有外部校准来源。接口的关闭不会以一次决定的形式发生，它会以产业结构的缓慢漂移出现：资源持续流向回收周期最短的方向，而支撑接口的那些活动一项一项地失去供给。要看清楚这个过程，需要把产业演化的轨迹摆出来，看操控的对象在历史上如何一层一层地转移。

\subsection{终极产业的崛起}

\subhead{产业演化的三级下沉}

产业史可以按操控对象来分段，而这条线索里有一个清楚的方向：每一轮革命都把操控的对象下沉一级，同时改变能量与信息的成本结构。工业革命操控的是宏观物质，棉纱、钢铁、化工原料、燃料是它的对象，它解决的问题是温饱与动力，方法是把自然过程搬进工厂，用热机和机械替代人力与畜力。这一轮的成就可以用物质指标衡量：钢产量、铁路里程、发电量、粮食单产。它把人类从物质的绝对稀缺中部分解脱出来，代价是化石能源的消耗和随之而来的排放。

信息革命把操控对象从物质换成比特：电子在电路里的状态、信号的编码方式、计算的步骤、信息的流转路径。这一轮的成本结构与上一轮不同，因为复制信息的边际成本趋近于零，而这一性质使信息产业能够在几十年里把单位成本压低几个数量级。它用信息指标衡量自己：算力、带宽、存储、连接数。它解决的问题是计算与传播，而这两个问题在过去受限于人脑的处理速率和物理介质的传递速度。

智能革命把操控对象推进到第三层：细胞、基因表达、蛋白质结构、代谢通路、神经活动，也就是生命过程本身。这一层与前两层有一个关键差别：它的对象不是惰性的，而是自己维持自身的系统，有历史、有反馈、会对外部干预作出反应。修改一行代码，得到的是修改后的代码；修改一个基因组，得到的是一个会回应的过程，而这个回应可能在几年之后才显现。这个差别决定了这一轮的产业形态不会重复前两轮的成本曲线，也不会重复前两轮的风险结构。

需要说明这个分段不是按重要性排序，也不是说前两轮已经结束。按本书的用法，这三层对应三种介质：原子、比特与生命，第一层操控的是原子构成的物质，第二层操控的是比特承载的信息，第三层操控的是会自己维持自身的生命过程。物质与信息仍然构成经济的主体，而且在很长时间里都会如此。分段的作用在于指出稀缺性转移的方向：物质与信息的稀缺性在被持续压缩，而第三层的约束（发育的周期、代谢的速率、免疫的应答、老化的进程）没有跟着一起下降。

\subhead{生命、心智与信任}

生命方向的轮廓已经可以看见：结构预测与分子设计把候选的生成成本压得极低，基因编辑把修改遗传信息变成可编程的操作，\cite{jinek2012programmable} 精准医疗把治疗方案与个体特征对齐，老年医学把重点从治疗急性病扩展到延缓功能衰退。这个方向的需求侧有很高的确定性：全球出生时预期寿命在 2024 年约为 73.3 岁，比 1995 年提高了 8.4 岁；65 岁以上人口已占全球人口约十分之一，2050 年前后将升到约六分之一，到 2070 年代末 65 岁以上人口将达到 22 亿，并超过 18 岁以下人口的总数。\cite{un2024wpp}

需求增长遇到的是一个供给速率受限的部门。William J. Baumol 与 William G. Bowen 在 1966 年对表演艺术的经济研究中，以及Baumol 1967 年的论文中，给出了一个后来被反复引用的解释：当制造业的劳动生产率因为技术而持续提高时，那些必须由人一对一完成的服务无法获得同样幅度的生产率增长，于是它们的单位成本相对于其他部门持续上升。\cite{baumol1966performing,baumol1967macro} 照护、陪伴、心理治疗、教育都属于这一类，它们的时间不能被压缩，一对一的比例不能被改变，因此它们的相对价格会长期上升。这个机制还掩盖了一个更基本的事实：这些活动的大部分从来没有进入定价，它们的成本一直由家庭内部的非支付劳动承担。

把这条机制写成相对成本的走向就是，

\begin{equation}
\frac{C_{\text{慢}}(t)}{C_{\text{快}}(t)}=\frac{C_{\text{慢}}(0)}{C_{\text{快}}(0)}\,e^{\,r t}
\end{equation}

式中 $C_{\text{慢}}$ 为生产率无法大幅提高的部门的单位成本，单位为货币每单位产出；$C_{\text{快}}$ 为生产率持续提高的部门的单位成本，单位相同；$r$ 为进步部门劳动生产率的增长率，单位为时间的倒数；$t$ 为时间。劳动力可以在部门之间流动时，两部门的工资同步上升，进步部门的单位成本被自身生产率的提高抵消，慢部门的单位成本则按 $r$ 上升，因此两者的比值按指数扩大。这个式子的含义是：慢部门变贵不是因为它们变得更重要，而是因为它们无法被加速，这与正文关于高黏滞位置的判断落在同一批对象上。适用边界有两处：$r$ 的数值须由分部门的统计数据确定，本式不给出也不需要给出它的取值；它只描述相对成本的走向，不构成对慢部门的评价，正文明确只取这条相对价格上升的机制，不取成本病理论对部门本身的判断。

AARP公共政策研究所在2026年3月的报告中估计，美国2024年成人家庭照护约有495亿小时，按所采用方法折算的经济价值约为1.01万亿美元。\cite{aarp2026valuing} 这是估值而非实际支付，且方法已有调整，不能把它与旧报告的差额全当作同口径增长。未付薪仍然意味着有人实际投入了时间和能力。

心智方向涉及心理健康、支持与关系。WHO于2025年9月更新的报告指出，全球超过十亿人生活在精神障碍之中，服务仍有显著缺口。\cite{who2025mental} 对话服务可以提供帮助，临床效果、安全、升级处置、隐私与可达性则应按其实际宣称分别建立证据。

信任方向处理的是判断的可依赖性。2024 年 1 月，一家跨国工程公司的香港办公室发生过一起案件：一名财务人员在一场视频会议中看到了他熟悉的同事与上司，会后按指示分 15 笔转账，共计约 2 亿港元，折合 2500 万美元左右，一周后与总部核实时才发现，会议中除他以外的所有与会者都是深度伪造。\cite{hk2024deepfake} 这件事的意义不在金额，而在于它关闭了一条长期被依赖的验证通道：视频中的脸和声音，从这一刻起不再构成身份证明。信任需要重新锚定在不易被低成本伪造的东西上：双人复核、肉身见面、可追溯的签署、能够承担赔付责任的机构。

物理质押、可核查记录与可执行赔付，提供了几种增强可信度的思路。它们各自要回答：质押如何匹配后果，记录支持什么判断，谁决定触发救济，资金能否兑现。验证工具与制度可以降低建立信任的成本；生成能力扩大，也可能增加需要验证的对象。净结果需要连同两端一起测量。

\subhead{负熵的定价与它的分配}

把三个方向放在一起，可以看到一个共同点：它们的价值都锚在不可回滚的代价上。生命方向锚在身体的一次性，个体的病史、衰老过程和被干预过的机体都只发生一次，无法回滚到干预之前；心智方向锚在时间的不可压缩，关系的形成没有加速路径，信任的积累速率由兑现承诺的速度决定；信任方向锚在承诺的不可撤销，一次背约不能被撤回，只能被后续行为缓慢稀释。这三样东西都不是能力问题，因此不随算力的增长而贬值。

由此可以把这一轮的产业转移说成一个定价问题：在生成成本趋近于零之后，稀缺的不再是信息，而是那些不能被伪造的代价。有人用负熵来命名这些服务的性质：它们的作用不是提供信息，而是维持一个低熵状态。一个人的身体、一段关系的稳定、一个判断的可依赖性，都需要持续投入才能不退回混乱。这个命名需要克制地使用，因为奢侈在这里描述的是稀缺性，不是炫耀性；把这些服务说成奢侈品，容易被读成对昂贵的辩护，而它们的昂贵来自供给速率，不来自品质等级。

分配问题紧跟着出现。如果照护、心智服务与信任中介的稀缺性持续上升，而定价由市场完成，那么最需要这些服务的人往往最负担不起。这个结果与前面关于高黏滞位置的分析一致：稀缺会吸引更高的价格，但价格不能补偿全部代价，也不能解决准入问题。因此这些方向的扩张不会以纯市场的方式完成，它必然伴随制度安排：公共支付、社会保险、家庭与志愿的分担、对服务标准的强制。把这一层漏掉，所谓终极产业就会变成一份只有一部分人买得起的服务清单。

\subhead{这个判断的边界}

三个方向都有被低成本仿造侵蚀的风险，而侵蚀的形式比技术替代更隐蔽。生命关怀的形态可以被表演：以探访次数计量的陪护，以话术替代处置的咨询，以流程替代判断的照护，都能通过指标检查而不提供实质内容。心智服务可以被包装成消费品，把体验做成可以购买的项目，而其中不包含任何长期关系。信任中介可以被形式化：盖章式的背书，可以购买的评级，不承担赔付责任的担保。判断的判据与前面几节相同，仍然是代价能不能被伪造。如果一个生命关怀服务在判断出错时不承担任何后果，它就退回到低黏滞区，无论它的名字多么贴近人的需要。

另一处边界关于职业预期。把三大方向说成终极产业，不等于说这些工作更好、更有意义或者更值得追求，也不等于说进入这些领域就能获得稳定报酬。它们之中的很大一部分目前报酬偏低，而这个错配的纠正依赖制度安排，不依赖个人的选择。这个判断能支持的只有一件事：如果价值确实向不可回滚的代价移动，那么这些方向的相对重要性会上升，而上升的幅度与速度取决于社会愿意为它们支付多少。

\subsection{终极预言：自组织盖亚}

\subhead{从制造逻辑到生命逻辑}

工程系统按图纸被制造，按手册被维护，按预算被替换。这套逻辑有一个隐含前提：系统的状态可以被回滚：出错就修，坏了就换，过时就淘汰。生命系统不满足这个前提，因为它的状态由自身的历史累积而成，而维持自身运转需要持续的耗散。两者的差别可以写得很具体，写成三个条件。

第一个条件是自复制：系统能够制造自身的复制装置，而不只是制造产品。冯·诺依曼在 1966 年给出了这个条件的构造性结果：一部能自我复制的自动机必须同时携带一份描述和一部能读描述并照着建造的构造器，两者缺一不可：只有描述就只剩停滞，只有构造器就只剩制造。\cite{vonneumann1966theory}

第二个条件是自维持：系统通过持续耗散把自身保持在远离平衡的状态。薛定谔在 1944 年用“以负熵为食”提出过这一点，普里戈金与 Gr\'egoire Nicolis 在 1977 年给出了它的动力学形式：开放系统在远离平衡处可以通过物质与能量的通量形成并维持有序结构，代价是外部熵产的持续增加。\cite{schrodinger1944,nicolis1977,prigogine1977nobel}

第三个条件是可变异与可继承：复制带有误差，而误差能够被传递下去。有了这一条，系统的工作参数由选择过程设定，而不由设计者设定。

复制、维持和可遗传变异，使系统能够在环境中进入选择过程。复制保真度影响序列能够延续的方式。设每个位置独立出错的概率为$\nu$，序列长度为$\lambda$，则完整复制概率为

\begin{equation}
Q=(1-\nu)^{\lambda}\approx e^{-\nu\lambda}\qquad(\nu\ll1)
\end{equation}

$Q$为整条序列无错误复制的概率，$\nu\lambda$为期望错误数。$\nu\lambda$接近1时，无错误副本减少，但后代仍可与原序列高度相似。选择能否维持特定序列，还取决于相对适合度、变异网络与种群过程。DNA聚合酶的选择、校对与错配修复共同提高复制保真度。\cite{kunkel2000dna,gilbert1986}

序列变长时，要保持相同的完整复制概率，需要相应降低每位错误率。生命通过修复、校对与组织结构维持这些条件，具体能耗和收益由机制决定。保真度的提高因此需要进入资源账本，但不存在每降一个数量级就必须增加一层机制的固定换算。

\subhead{盖亚判据}

把尺度放大到行星，会得到一个更强的问题：如果一个星球上的工程系统与生物系统耦合到足够深的程度，整体是否也满足上面三个条件？

把自维持放到行星尺度，会出现一个更强的问题：工程网络能否在持续能源输入下，提高修复、循环与再生产其物质基础的能力？“盖亚”在这里组织这个问题。已有工业链可以复制许多设备，仍依赖人员、全球分工与特定原料，不能据此宣布整个系统已能从原始资源独立复制自身。

自维持、再生产与可继承变化需要逐项界定。某些状态不可撤销，会增加治理难度，却不是自组织的充分判据。对确定材料、边界与观察期，可以先记录一种较窄的量：循环返回可用供应的比例。

\begin{equation}
\kappa=\frac{\dot m_{\text{循环}}}{\dot m_{\text{总}}}
\end{equation}

分子计入本期返回并实际进入可用供应的材料，分母计入同一边界的全部可用投入，前者应是后者的一部分，比例才在0与1之间。库存、循环时滞、品质变化和重复计数都需说明。高循环比例可以减少新开采需求，不能独自证明能源自足、完全再生产或生态稳定。

循环比例应按具体材料、品质、边界与年份比较。钢、铝、建筑骨料与复杂复合材料的回收条件差别很大；回收率、再生投入比例与闭环保持品质的比例也不是同一个数。\cite{miatto2017how} $\kappa$能说明所定义的材料流，评价整个系统的自维持还须检查能源、损耗、维修和再生产依赖。

\subhead{这个预言的边界}

这个预言有三处边界必须写清楚，否则它会退化成一份技术清单之外的另一种东西。

第一，它目前不可检验。上面给出的是形式条件，不是时间表；说一个系统在结构上可能满足某个条件，与说它将会满足，是两件不同的事。可检验的部分只有那些能被计量的量（复制保真度、循环闭合度、单位时间的熵产），而这些量的历史序列都很短。

第二，它不含目的。盖亚在流行叙述里常被读成地球会照顾好自己，这是一种目的论读法，把“系统会维持某个状态”与“系统会朝着对某一方有利的方向走”混为一谈。形式上能支持的只有前者，而且它的维持方式未必对任何一方有利：一个自维持的行星系统完全可以在没有人的情况下运转。把维持本身读成善意，是把描述的结论当成规范的前提用。

\key{自复制的信息容量由复制保真度与基因组长度共同限定；工程系统一旦越过自维持的门槛，就不再是可以被单方面关停的对象。}第三处边界正是这条结论：关停一个能自我修复、并且已经在复制自身的系统，与关停一条产线，成本结构完全不同：前者要面对的是重新生长出来的部分，后者只需要支付停工的代价。这条差别落在代价能否被撤回上，它是本章四个原则在物理一侧的依据。

把本章的结论与四个原则对齐：阻尼本体论说明为什么代价不可替代，不完备性原则说明为什么系统必须保留外部的纠正通道，共生原则说明为什么分工不能被读成竞争，递归优化原则说明为什么调整必须一层一层地做下去。四条原则合起来不规定具体做什么，它们规定的是当不同的价值发生冲突时按什么次序裁决。

这部宪法有四块基石。第一块是保留不完美的权利：任何一个封闭的系统都不能在自身内部证明自己的目标是对的，因此必须保留可以被外部质疑、申诉和回滚的位置；把一个系统做到无懈可击，等于把纠正它的通道一并关上。第二块是定义代价的权力：谁承担代价，谁就有权参与定义什么算是代价，以及什么代价不能被转嫁。第三块是保留摩擦的分工：两侧的长处不同，联合才有内容，因此哪些阻力必须保留不能由快的一侧单独决定，分工也不能被读成竞争。第四块是定向：递归优化最终应当落在生命本身，落在人的身体、关系、判断和有限性上，而不是落在产出指标上。这四块基石之间没有高低之分，它们互为条件，缺掉任何一块，其余几块都会失去约束力。

“硅碳共生”描述一种互补安排：利用可复制的能力扩大选择，同时把行动接回证据、身体、权利与长期责任。它不是按材料永久划定分工。数字系统同样是物质的，生命也能学习、计算和委托；具体接口可以随证据与技术继续改变。

\key{文明把资源投向哪里，最终取决于它认为什么代价不可替代。}这个问题在生成能力充足之后不会消失，只会变得更清楚：向内的舒适与向外的高阻力前沿用的是同一笔资源，而两种选择给出的是两种不同的文明。前者并不缺少繁荣，它缺少的是与真实世界接触的接口。这个前景在下一章会以一种极端的形式出现：当一个系统获得了自我复制的物理能力，它就从制造逻辑进入了生命逻辑，而生命逻辑不接受任何一种单方面的收场方式。

\furtherreading{%
William J. Baumol and William G. Bowen，\textit{Performing Arts: The Economic Dilemma}\\
Robert H. Frank，\textit{Passions Within Reason: The Strategic Role of the Emotions}\\
James C. Scott，\textit{Seeing Like a State}\\
Karl Polanyi，\textit{The Great Transformation}\\
Zygmunt Bauman，\textit{Liquid Modernity}，Polity Press，2000\\
Elinor Ostrom，\textit{Governing the Commons}\\
Hans Jonas，\textit{The Imperative of Responsibility}\\
Amartya Sen，\textit{Collective Choice and Social Welfare}}


% <<<<<<<<<< parts/book-ch06.tex <<<<<<<<<<


% >>>>>>>>>> parts/book-ch07.tex >>>>>>>>>>
%% ===================== 本章参考文献 =====================
\printbibliography[heading=subbibliography,title=参考文献]
